\textbf{VI -- LA COLONNA APERTA}

Ricostruendo le partite proposte il lettore si sarà già reso conto della
importanza di occupare una colonna aperta. Essa è come un'autostrada che
porta dritto in territorio nemico.

In concreto il vantaggio del controllo di una colonna aperta (per
semiaperta, abbiamo già avuto occasione di dire che si intende una
colonna dove è posto un pedone di colore contrario) risiede in più
fattori. Innanzitutto c'è il vantaggio generico di avere a disposizione
più spazio, dove più spazio si traduce in più mosse possibili ed è noto
che più mosse si hanno a disposizione più se ne hanno di buone. Ma il
vero vantaggio risiede nella possibilità di trasformare il controllo
verticale in orizzontale con l'occupazione della settima traversa, dove
in genere è allineato il maggior numero di pedoni, o dell'ottava, la
traversa del Re.

\textbf{a) La settima traversa}

Si consideri questa posizione:

\bookdiagramplaceholder{assets/diagrams/image89.png}{89}

Il vantaggio del Bianco è chiaro. La Torre in settima esercita una
pressione su tutto lo schieramento nero: non può essere scacciata dal Re
e inchioda la Torre avversaria alla difesa del Pb7.

La partita che segue è un classico esempio dei danni che una Torre in
settima può provocare in centro partita.

Steinitz--Von Bardeleben, Hastings 1895 \textbf{(Partita Italiana)}

\textbf{1.e4 e5 2.¤f3 ¤c6 3.¥c4 ¥c5 4.c3 ¤f6 5.d4 exd4 6.¤xd4 ¥b4+
7.¤c3} l'idea di questa mossa è di guadagnare tempi col sacrificio due
pedoni. I pericoli cui il Nero va incontro in caso di accettazione erano
già stati indicati da Gioacchino Greco nel XVII secolo. Per esempio:
7... ¤xe4 8.0-0 ¤xc3 9.bxc3 ¥xc3 (9... d5!) 10.¥a3! (la mossa di Atkins)
10... d5 11.¥b5 ¥xa1 12.¦e1+ ¥e6 13.¤e5 e il Nero non ha difesa. Nel
nostro secolo schiere di analisti hanno analizzato a fondo questa
continuazione giungendo alla conclusione che le sue estreme conseguenze
non sono sfavorevoli al Nero; per questa ragione oggi si preferisce
giocare 7.¥d2 col possibile seguito 7... ¥xd2+ (7... ¤xe4 8.¥xb4 ¤xb4
9.¥xf7+! ¢xf7 10.£d3+ d5 11.£xb4) 8.¤bxd2 d5!

\textbf{7... d5} oggi si preferisca giocare 7... ¤xe4 ma in mancanza di
analisi approfondite la scelta del Nero è del tutto sana e giustificata.

\textbf{8.exd5 ¤xd5} il Bianco puntava ad avere un forte centro e si
ritrova con un pedone debole.

\textbf{9.0-0 ¥e6} non si può accettare il pedone: 9... ¤xc3? 10.bxc3
¥xc3 11.¥a3! ¥b4 12.¦e1+ ¢f8 13.¥xb4+ ¤xb4 14.£d3

\textbf{10.¥g5 ¥e7 11.¥xd5 ¥xd5 12.¤xd5 £xd5 13.¥xe7 ¤xe7 14.¦e1}
vantaggi e svantaggi si equivalgono. Il Bianco ha un pedone debole in
d4, il Nero ha il problema di trovare un rifugio al Re.

\textbf{14... f6} preparando \emph{l'arrocco artificiale} (¢f7, ¦e8,
¢g8).

\textbf{15.£e2 £d7 16. ¦ac1 c6?} si doveva giocare 16... ¢f7. Ora il
Nero non riesce più a completare lo sviluppo.

\textbf{17.d5!} sgombera la casa d4 al Cavallo.

\textbf{17... ¤xd5 18.¤d4} minaccia il salto in f5.

\textbf{18... ¢f7 19.¤e6 ¦hc8 20.£g4 g6} è sbagliata 20... ¦xc1?
21.£xg7+ ¢e8 22.£xf8 matto.

\textbf{21.¤g5+ ¢e8}

\bookdiagramplaceholder{assets/diagrams/image90.png}{90}

\textbf{22.¦xe7+!!} una mossa molto bella. Ora non va 22... ¢xe7 per
23.¦e1+ e ora:

\emph{a)} 23... ¢d6 24.£d4+ ¦c5 \emph{(24... ¢c7 25.¤e6+ ¢b8 26.£f4+)}
25.¦e6+;

\emph{b)} 23... ¢d8 24.¤e6+ ¢e8 25.¤c5+ ma il Nero dispone di una
replica ingegnosa.

\textbf{22... ¢f8} il Nero contrappone alla Torre in settima del Bianco
la debolezza della prima traversa bianca! La Donna è difesa grazie alla
minaccia di matto in c1 e inoltre tutti i pezzi dell'avversario si
trovano sotto minaccia di presa, ma Steinitz ha visto più lontano!

\textbf{23.¦f7+! ¢g8} impossibile giocare 23... £xf7 per 24. ¦c8+ mentre
a 23... ¢e8 segue 24.¦e1+.

\textbf{24.¦g7+!} nella settima traversa questa Torre fa il bello e il
cattivo tempo.

\textbf{24... ¢h8} non salva 24... ¢f8 per 25.¤xh7+ ¢xg7 (25... ¢e8
26.¦e1+) 26.£xd7+.

\textbf{25.¦xh7+!} una vera persecuzione!

\textbf{25...} il Nero abbandona

Fu lo stesso Steinitz a fornire, al termine della partita, la precisa
dimostrazione di vittoria. 25... ¢g8 26.¦g7+ ¢h8 27.£h4+ ¢xg7 28.£h7+
¢f8 29.£h8+ ¢e7 30.£g7+ ¢e8 \emph{(30... ¢d6 31.£f6+; 30... ¢d8
31.£f8+)} 31.£g8+ ¢e7 32.£f7+ ¢d8 33.£f8+ £e8 34.¤f7+ ¢d7 35.£d6 matto.

Una partita memorabile!

Dopo una partita così complicata eccone un'altra di una semplicità
disarmante.

La semplicità era la peculiarità di Capablanca che sembrava vincere le
partite facendo uso di mosse lapalissiane. Nella partita che segue non
sembra accadere nulla, i cambi si succedono ai cambi, eppure il Bianco è
costretto alla resa dopo appena 23 mosse.

Presento la partita senza commenti; il lettore vi riconoscerà, oltre al
tema in esame (colonna aperta e sfruttamento della settima traversa),
anche un tentativo di attacco di minoranza del Bianco (15.a4 e 19.¦fb1)
frustrato dalla semplice mossa di Capablanca 20... a6 (21.a5 b5!).

Alatorzev--Capablanca, Mosca 1935 \textbf{(Gambetto di Donna Rifiutato)}

1.d4 ¤f6 2.c4 e6 3.¤c3 d5 4.¥g5 ¥e7 5.e3 0-0 6.¤xd5 ¤xd5 7.¥xe7 £xe7
8.¤f3 ¤xc3 9.bxc3 b6 10.¥e2 ¥b7 11.0-0 c5 12.¤e5 ¤c6 13.¤xc6 ¥xc6 14.¥f3
¦ac8 15.a4 ¤xd4 16.¤xd4 g6 17.¥xc6 ¦xc6 18.£d3 £d7 19.¦fb1 ¦fc8 20.h3 a6
21.£a3 ¦c2 22.£d6 ¦xf2 23.£g3 ¦e2 0--1

\textbf{b) Debolezza dell'ottava traversa}

Nella fase finale della partita Bernstein--Capablanca abbiamo già avuto
modo di renderci conto a quali tatticismi può portare un'ottava traversa
non sufficientemente protetta quando il Re sia privo di case di fuga in
settima.

\bookdiagramplaceholder{assets/diagrams/image91.png}{91}

La posizione del diagramma illustra a dovere lo sfruttamento dell'ottava
traversa. Essa è tratta da una partita attribuita a Adams e Torre (New
Orleans 1920) ma, in realtà si tratta di una costruzione a tavolino
ideata da Torre come regalo al suo amico Adams.

Ciò non toglie che rimanga un esempio di spettacolare efficacia.

\textbf{18.£g4!} inizio di una caccia alla Donna. È evidente che, per
via della debolezza dell'ottava traversa, la Donna nera non può lasciare
la diagonale a4-e8.

\textbf{18... £d5 19.£c4!} una mossa molto bella. La Donna bianca
continua a mettersi impunemente in presa.

\textbf{19... £d7 20. £c7!} entrando nel pollaio!

\textbf{20... £d5 21.a4} notare che in caso di 21.£xb7 è il Nero a
sfruttare la debolezza dell'ottava (prima) traversa! 21... £xe2! 22.¦xe2
¦c1+ e matto alla seguente.

\textbf{21... £xa4 22.¦e4} la debolezza dell'ottava ha irrigidito i
pezzi neri e il Bianco può fare i suoi comodi.

\textbf{22... £d5 23.£xb7} tutte le case della diagonale sono ora
gua¢date dal Bianco. Il Nero abbandona.

Per finire una partita del grande Alekhine che, consolidata la conquista
di una colonna spinge addosso all'arrocco avversario i pedoni per
indebolire l'ottava traversa.

Evenson--Alekhine, Kiev 1916 \textbf{(Difesa Philidor)}

\textbf{1.e4 e5 2.¤f3 d6 3.d4 ¤f6} una mossa introdotta da Nimzowitsch.
In un periodo in cui i sostenitori della Scuola Classica si
accapigliavano con gli Ipermoderni, Alekhine si tenne fuori dalle
dispute teoriche ma, con vero spirito pragmatico, senza alcun
pregiudizio prendeva quel che c'era di buono in entrambe le scuole. Non
solo, quindi, adottò varianti indicate da Nimzowitsch ma in seguito sarà
lui stesso, fatto tesoro di quanto di buono era ravvisabile nelle nuove
idee, a introdurre nella pratica una nuova apertura di stampo
ipermoderno (1.e4 ¤f6).

\textbf{4.¤c3} oppure 5.£xe5 ¤xe4

\textbf{4... ¤bd7 5.¥c4 ¥e7 6.0-0} interessante, ma probabilmente non
del tutto corretto, è il sacrificio 6.¥xf7+ ¢xf7 7.¤g5+ ¢g8 8.¤e6 £e8
9.¤xc7 £g6 10.£xa8 £xg2.

\textbf{6... 0-0 7.£xe5 £xe5 8.¥g5 c6} evita eventuali salti in d5 (e
b5) e prepara la casa c7 alla Donna. Incidentalmente minaccia anche
l'espansione sul lato di Donna con b7-b5, perciò la prossima mossa del
Bianco è quasi forzata.

\textbf{9.a4 £c7 10.£e2 ¤c5 11.¤e1} una mossa involuta. Sembra migliore
11.h3 e se 11... ¤e6 12.¥e3 ¦d8 13.¦fd1. Così il Nero conquista la casa
d4.

\textbf{11... ¤e6 12.¥e3 ¤d4 13.£d1 ¦d8} Alekhine si sviluppa a suon di
minacce. Evenson ha problemi a coordinare i pezzi e già ora si può
considerare il Nero in vantaggio.

\textbf{14.¤d3 ¥e6 15.¥xe6 ¤xe6 16.£e1 ¦d7} in una situazione in cui il
Bianco ha problemi nel cootdinare l'azione dei pezzi, il Nero ha tutto
il tempo per ra£doppiare le Torri sulla colonna d.

\textbf{17.f3} non va 17.f4 per 17... ¤g4. In linea generale è sbagliato
aprire le linee quando si è in arretrato di sviluppo o si hanno disposti
in modo meno armonico.

\textbf{17... ¦ad8} il Nero ha raddoppiato le Torri sulla colonna ``d''
ma la presenza del ¤d3 ne impedisce lo sfruttamento. Con la prossima
manovra di Cavallo il Nero tenderà alla rimozione di questo ostacolo.

\textbf{18.¥f2 ch5 19.¤e2 c5 20.b3 chf4 21.¤exf4 ¤xf4 22.¤xf4 exf4 23.c3
£e5 24.¦a2 ¦d3 25.¦c2 b6 26.£c1 £e6 27.£d1 ¥f6 28.b4 c4 29.£c1} il
problema del Bianco è di non poter contrastare in alcun modo, al Nero,
il controllo della colonna d.

\textbf{29... g5}

\bookdiagramplaceholder{assets/diagrams/image92.png}{92}

Le spinte a fondo dei pedoni sul lato di Re stringeranno il Bianco in
una morsa e preludono allo sfruttamento della colonna d.

\textbf{30.h3 ¥e5 31.£a1 h5 32.a5 g4 33.axb6 axb6 34.¥h4 f6 35.¥e1 g3}
Alekhine sta letteralmente preparando la bara a Bianco. La spinta in
profondità di questo pedone prepara la penetrazione delle Torri nella
prima traversa, divenuta terribilmente debole. Come se non bastasse,
l'attività dei pezzi bianchi è ridotta al lumicino e non consente alcuna
controffensiva.

\textbf{36.£a6 £c6 37.£a3 b5 38.£d2 £d6+ 39.¢h1 ¦d1} lo scopo è
semplice: 40... £e3 41.¦e2 £xc3. La posizione del Bianco si fa sempre
più difficile.

\textbf{40.¦c1 £e3 41.¦a1 ¥c7 42.£a2 ¦xa1 43.£xa1 £e2 44.¢g1 ¥b6} il
matto (45... ¥xg1 46.¢xg1 £xe1 47.£xe1 £xe1 matto) è imparabile. Il
Bianco abbandona.

\textbf{c) Il punto d'appoggio}

\bookdiagramplaceholder{assets/diagrams/image93.png}{93}

chi ha il tratto, supponiamo il Bianco, può, sfruttando il ``punto
d'appoggio'' in f5, forzare la conquista della colonna o, a scelta
dell'avversario, la formazione di un pedone passato. La semplice manovra
è 1.¦f5. Se il Nero cambia le Torri il Bianco si trova con un pedone
passato in f5 (è ovvio che il Bianco deve stabilire prima se il pedone
passato sia forte o debole), se il Nero non cambia il Bianco alla mossa
successiva gioca 2.¦af1 e conquista una colonna aperta.

In altri termini il punto d'appoggio è un caso particolare di casa,
debole nello schieramento nemico, che viene sfruttata per ra£doppiare le
Torri su una colonna aperta.
